\documentclass{beamer}
\usepackage{JEFtheme}
\usepackage{caption}
\usepackage{csvsimple}
\usepackage{expl3} 
\usepackage{etoolbox}
\usepackage{xstring}
\usepackage[german]{babel}
\newcommand\Committee{SEDE}
\newcommand\Fraktion{LEER}

\newcommand\Thema{Armee}
\newcommand\city{Ulm}
\newcommand\datum{11.05.2026}
\newcommand\timeVorb{07:30-09:00}
\newcommand\timeEinf{09:00-09:30}
\newcommand\timeFrakOne{09:30-11:00}
\newcommand\timePauseOne{11:00-11:15}
\newcommand\timeAuss{11:15-12:30}
\newcommand\timeMittag{12:30-13:00}
\newcommand\timeFrakTwo{13:00-13:30}
\newcommand\timeFinVerh{13:30-14:00}
\newcommand\timePlenar{14:00-15:15}
\newcommand\timeDebr{15:15-15:30}
\newcommand\politiker{Andrea Wechsler}
\newcommand\politikerOffice{Mitglied des Europäischen Parlaments}
\newcommand\stadtvertreter{Nalan Schmidt}
\newcommand\stadtvertreterOffice{Anlauf- und Koordinierungsstelle Bürgerdialog (Open Government) der Stadt Ulm}
\newcommand\localSupport{das Europe Direct Ulm}
\newcommand\sponsor{die Vertretung der Europäischen Kommission in München}
\newcommand\jefvorsitz{Farras Fathi}
\newcommand\gendervorsitz{Landesvorsitzender}
\newcommand\evpLeader{Tobias}
\newcommand\evpRoom{Kleiner Sitzungssaal}
\newcommand\sdLeader{Antonia}
\newcommand\sdRoom{Veranstaltungsraum Europe Direct}
\newcommand\reLeader{Alica}
\newcommand\reRoom{Besprechungszimmer Donaubüro I}
\newcommand\fifthLeader{Farras}
\newcommand\fifthRoom{Besprechungszimmer Öffentlichkeitsarbeit}
\newcommand\pfeLeader{Hannah}
\newcommand\pfeRoom{Besprechungszimmer BM III}
\newcommand\numSuS{27}
\newcommand\location{im Ulmer Rathaus}
\newcommand\fifthGroup{Linke}
\newcommand\anzahlcomm{3}
\newcommand\totcoms{26}
\newcommand\minmems{25}
\newcommand\maxmems{90}


% -------------------------------------------------------------

\IfEqCase{\Fraktion}{%
	{EVP}{
		\newcommand{\Fraktionsname}{Europäische Volkspartei}
	    \newcommand{\Fraktionskuerzel}{EVP}
	}%
	{SD}{
		\newcommand{\Fraktionsname}{Progressive Allianz der Sozialdemokraten}
        \newcommand{\Fraktionskuerzel}{S\&D}
	}%
	{RE}{
		\newcommand{\Fraktionsname}{Renew Europe}
		\newcommand{\Fraktionskuerzel}{Renew}
	}%
	{Green}{
		\newcommand{\Fraktionsname}{Die Grünen/Europäische Freie Allianz}
		\newcommand{\Fraktionskuerzel}{Grüne/EFA}
	}%
	{PfE}{
		\newcommand{\Fraktionsname}{Patrioten für Europa}
		\newcommand{\Fraktionskuerzel}{PfE}
	}%
	{Left}{
		\newcommand{\Fraktionsname}{Die Linke}
		\newcommand{\Fraktionskuerzel}{GUE/NGL}
	}%
}


%Information to be included in the title page:
\title{\Fraktionsname\ (\Fraktionskuerzel)}
\subtitle{2. Fraktionssitzung}
\date{\datum}

\begin{document}

\frame{\titlepage}

\begin{frame}{Ausschusssitzungen} % Besprechung Ausschusssitzungen
\vspace{-1.5cm}
Berichtet aus jedem Ausschuss:
\newline
\begin{itemize}
    \item Habt ihr euren Änderungsantrag durchbekommen?
    \item Welche anderen Fraktionen haben ihre Änderungsanträge durchbekommen?
    \item Wer hat mit wem zusammengearbeitet?
    \item Ist der Gesetzesentwurf nun mehr in unserem Sinne oder weniger?
\end{itemize}
\end{frame}

\begin{frame}{Finaler Änderungsantrag} % Festlegung finaler ÄA
\vspace{-0.5cm}
In der Plenarsitzung können wir als Fraktion noch einen (!) finalen Änderungsantrag stellen. Beispielsweise: \newline
\begin{itemize}
    \item Alter Änderungsantrag, der im Ausschuss abgelehnt wurde
    \item Verworfene Idee aus erster Fraktionssitzung
    \item Rückgängigmachen von Änderungsantrag, der uns überhaupt nicht passt
\end{itemize}
\end{frame}

\begin{frame}{Wahl des Fraktionsvorstands}
\vspace{-1.5cm}
\begin{itemize}
    \item Fraktionsvorsitzende:r hält Abschlussrede (ca. 1 Min.)
    \item Vize begründet finalen Änderungsantrag (ca. 30 Sek.) \& nimmt Stellung zu den ÄA der anderen Fraktionen \newline (jeweils ca. 10 Sek.)
\end{itemize} 
\end{frame}

\begin{frame}{Einteilung Verhandlungsgruppen}
%\vspace{-1.5cm}
\begin{itemize}
    \item Der Rest teilt sich einer der anderen Fraktionen zu, sodass sich vier gleich große Gruppen ergeben.
    \item Eure Aufgabe ist es mit den Verhandlungsgruppen der Fraktion, der ihr zugeordnet seid, über letzte Kompromisse zu diskutieren.
    \item Geht dabei ähnlich zur Ausschusssitzung vor, endgültige Änderungen an euren Änderungsanträgen (bspw. Kompromissformulierungen oder Zusammenlegungen von Anträgen) müsst ihr dem Parlamentspräsidium bis XXX melden.
\end{itemize} 
\end{frame}


\begin{frame}{Rede Fraktionsvorsitzende}
\vspace{-1.5cm}
Falls die beiden Fraktionsvorsitzenden noch nicht wissen, was sie sagen sollen, helfen wir euch gerne!
\end{frame}


\end{document}